\section{Introduction}\label{sec:introduction}

The H\'enon--Lane--Emden system is
\begin{equation}\label{eq:system}
 -\Delta u=|x|^A v^p,\qquad
 -\Delta v=|x|^B u^q
 \qquad\text{in }\R^n\setminus\{0\}.
\end{equation}
Its subcritical nonexistence conjecture asks whether the weighted
hyperbola alone rules out positive solutions, without symmetry or
conditions at infinity. The unweighted case \(A=B=0\) is the classical
Lane--Emden system.

The scalar subcritical Liouville theorem of Gidas and Spruck is a
starting point for this question~\cite{GidasSpruck1981}. For systems,
Mitidieri developed a Rellich identity and early nonexistence
regions~\cite{Mitidieri1993,Mitidieri1996}. Serrin and Zou proved the
unweighted three-dimensional assertion under a growth condition;
Pol\'a\v{c}ik, Quittner, and Souplet removed that condition, and
Souplet settled dimension four~\cite{SerrinZou1996,
PolacikQuittnerSouplet2007,Souplet2009}. Li, Li, and Wei obtained a
further unweighted nonexistence region for \(p,q\ge1\) in every dimension
\(n\ge5\)~\cite[Theorem~1.1]{LiLiWei2025}. Liouville theorems also serve
as inputs to rescaling arguments for a priori bounds and singularity
estimates~\cite{GidasSpruckBounds1981,PolacikQuittnerSouplet2007}.

For the weighted system, Bidaut-V\'eron and Giacomini identified the
critical hyperbola in the regular radial theory under admissible
weights~\cite{BidautVeronGiacomini2010}. Phan formulated the unrestricted
weighted question for \(A,B>-2\) as Conjecture~C and established
bounded, low-dimensional, and conditional cases~\cite{Phan}.
Fazly and Ghoussoub obtained bounded and stability-based results for
nonnegative weights~\cite{FazlyGhoussoub2014}; Li and Zhang subsequently
proved the weighted conjecture in dimension three without a boundedness
assumption~\cite{LiZhang2019}. Results in restricted exponent ranges
extend to dimensions four and five~\cite{Li2022}. More recently,
Huang and Zou stated the all-real-weight, origin-continuous assertion
as Conjecture~B and proved cases with stability, decay, or energy
hypotheses~\cite{HuangZou}. Xu and Luo obtained a higher-order weighted
theorem over the full exponent range for radial solutions subject to
regularity and growth conditions~\cite{XuLuo2026}. The full nonradial
question in higher dimensions is the subject of this paper.

\begin{theorem}\label{thm:main}
Let \(n\ge2\) be an integer, let \(p,q>0\), and let \(A,B\in\R\).
Suppose that
\begin{equation}\label{eq:subcritical}
 \frac{n+A}{p+1}+\frac{n+B}{q+1}>n-2.
\end{equation}
There are no functions
\[
 u,v\in C^2(\R^n\setminus\{0\})\cap C(\R^n),
 \qquad u(x)>0,\quad v(x)>0\quad(x\in\R^n),
\]
satisfying Equation~\eqref{eq:system} at every \(x\ne0\).
\end{theorem}

\begin{corollary}[Unweighted Lane--Emden conjecture]
\label{cor:lane-emden}
For every integer \(n\ge3\) and \(p,q>0\) satisfying
\[
 \frac{1}{p+1}+\frac{1}{q+1}>\frac{n-2}{n},
\]
there is no strictly positive \(C^2(\R^n)\) solution of
\[
 -\Delta u=v^p,\qquad -\Delta v=u^q\qquad\text{in }\R^n.
\]
\end{corollary}

In particular, positivity at the origin is part of the hypothesis,
but differentiability there is not. No radial symmetry, boundedness,
integrability, finite energy, stability, or growth or decay condition
at infinity is required.

Theorem~\ref{thm:main} resolves Huang and Zou's Conjecture~B positively
and yields Corollary~\ref{cor:lane-emden} by setting \(A=B=0\).
For \(n\ge3\) and \(A,B>-2\), its nonexistence conclusion joins the
radial existence theorem of Bidaut-V\'eron and
Giacomini~\cite[Theorem~1.4(i)]{BidautVeronGiacomini2010} to prove Phan's
Conjecture~C in full, including existence on the critical hyperbola.
Corollary~\ref{cor:phan-classification} records this derived classification.
The weak regularity at the origin is genuine: positivity and continuity,
rather than a prescribed order of vanishing or smoothness there, are all
that is assumed.

The weighted assertion cannot be inferred merely by applying the
unweighted assertion at the same exponents. For instance,
\(n=5\), \(p=q=4\), \(A=B=3\) satisfies
Equation~\eqref{eq:subcritical}, whereas
\(1/(p+1)+1/(q+1)=2/5<3/5\).
Likewise, the unweighted input in Phan's transfer argument must remain
within its actual scope; weighted subcriticality does not imply that
input~\cite{Phan}.

\paragraph{The weighted argument.}
Rellich--Pohozaev identities and spherical estimates have long been
central to system Liouville proofs~\cite{Mitidieri1993,Souplet2009,Phan}.
Chen and Li used Newtonian integral representations with moving planes
under integrability hypotheses~\cite{ChenLi2009}; Cheng and Huang later
formulated an energy-estimate criterion for the unweighted
conjecture~\cite{ChengHuang2019}. The passage from a uniform local
weighted-energy bound to nonexistence by origin-centered dilation also
appears in Li~\cite[Section~2]{Li2022}. Here a localized potential-kernel
virial identity and an energy-relative interval-pressure estimate supply
that bound throughout the subcritical range.

There are two difficulties in passing from unweighted source layers
to Equation~\eqref{eq:system}. Translations do not preserve the
equation, and an endpoint of a layer
\(\{|x|^B u(x)^q>t\}\) has value
\((t/|x|^B)^{1/q}\), not the constant \(t^{1/q}\).
We address both directly.

First, origin-centered dilations give uniform source masses and
truncated-potential bounds on a fixed ball. A compactly supported
virial identity then retains the exact spatial-weight coefficients.
Second, the interval comparison is applied with directions in a
small cap around the radial direction and with weights constant on
each retained interval. Retained intervals are short enough that
their endpoint weight variation is small. The discarded intervals
are controlled by polar integration of source masses and a
high/low layer decomposition.

All errors are arbitrarily small relative to a fixed localized
energy, up to a universal additive constant. No bound on that energy
is assumed in obtaining the error estimates. The weighted
subcritical gap then supplies such a bound, and dilation gives a
contradiction. This combination of spatially varying directional
cutoffs, intervalwise infima, and energy-relative layer errors is
the part of the method that may be useful for other weighted
potential systems. The elementary interval comparison itself is
reproved below.

\paragraph{Organization.}
Section~\ref{sec:preliminaries} treats the origin, the low-exponent
cases, and the potential estimates.
Section~\ref{sec:localization} proves the localized virial identity
and compares the two localized interaction marginals.
Section~\ref{sec:intervals} gives the geometric interval inequality.
Section~\ref{sec:pressure} proves its weighted pressure consequence,
including the needed layer-error bounds.
Section~\ref{sec:conclusion} completes the proof by scaling and derives
the existence--nonexistence classification from the radial theorem.
